% Template de artigo curto para o MÁGICA — SBGames 2026.
% Projeto: The King is Dead.
% Compilador recomendado no Overleaf: pdfLaTeX.
% Antes de submeter, remova todos os comandos \preencher{...}.

\documentclass[12pt]{article}

\usepackage{sbc-template}
\usepackage[brazilian]{babel}
\usepackage[T1]{fontenc}
\usepackage[utf8]{inputenc}
\usepackage{graphicx}
\usepackage{booktabs}
\usepackage{tabularx}
\usepackage{array}
\usepackage{xcolor}
\usepackage{fancyhdr}
\usepackage{microtype}
\usepackage[hidelinks]{hyperref}

\definecolor{instrucao}{RGB}{145,35,35}
\newcommand{\preencher}[1]{\textcolor{instrucao}{\textbf{[Preencher:} #1\textbf{]}}}
\newcolumntype{Y}{>{\raggedright\arraybackslash}X}

% Corrige o rótulo da versão original do estilo SBC.
\renewcommand{\palavraschave}[1]{%
  \vspace{6pt}\small\linebreak
  \textbf{\textit{Palavras-chave---}} #1
}

% Cabeçalho solicitado pelo MÁGICA no padrão dos artigos da SBC.
\fancypagestyle{magica}{%
  \fancyhf{}
  \fancyhead[C]{\footnotesize
    XXV Simpósio Brasileiro de Jogos e Entretenimento Digital (SBGames 2026) -- Goiânia/GO\\
    \textbf{Atividade: MÁGICA -- Workshop de Games na Graduação e na Educação Básica}}
  \renewcommand{\headrulewidth}{0pt}
  \renewcommand{\footrulewidth}{0pt}
}
\setlength{\headheight}{28pt}

\title{The King is Dead: desenvolvimento incremental de um jogo de estratégia em Python\\
\medskip
\small{\textit{Title: The King is Dead: Incremental Development of a Strategy Game in Python}}}

% O(s) estudante(s) deve(m) aparecer primeiro; o professor responsável, por último.
\author{Nome do Estudante\inst{1}, Nome do Estudante\inst{1}, Nome do Professor\inst{1}}

\address{Nome do Curso -- Nome da Instituição\\
  Cidade -- UF -- Brasil
  \email{estudante@exemplo.edu.br, professor@exemplo.edu.br}
}

\begin{document}

\maketitle
\thispagestyle{magica}
\pagestyle{empty}

\begin{abstract}
\textbf{Introduction:} \preencher{contextualize o projeto e o problema em uma ou duas frases.}
\textbf{Objective:} \preencher{declare o objetivo do jogo e deste relato.}
\textbf{Methodology or Steps:} \preencher{resuma o desenvolvimento incremental, a arquitetura e as entregas.}
\textbf{Results:} \preencher{informe o que está funcional e as evidências de validação.}
\end{abstract}

\keywords{digital games, game development, state tree, Python, Pygame.}

\begin{resumo}
\textbf{Introdução:} \preencher{contextualize o projeto e o problema em uma ou duas frases.}
\textbf{Objetivo:} \preencher{declare o objetivo do jogo e deste relato.}
\textbf{Metodologia ou Etapas:} \preencher{resuma o desenvolvimento incremental, a arquitetura e as entregas.}
\textbf{Resultados:} \preencher{informe o que está funcional e as evidências de validação.}
\end{resumo}

\palavraschave{jogos digitais, desenvolvimento de jogos, árvore de estados, Python, Pygame.}

\section{Introdução}

\preencher{Apresente o contexto acadêmico, a disciplina ou projeto, o gênero do jogo e o problema enfrentado. Explique por que uma adaptação de \textit{The King is Dead} é relevante para exercitar modelagem de regras, interfaces e inteligência artificial. Finalize indicando a organização do texto.}

% Se necessário, cite o manual do jogo: \cite{osprey2020king}.

\section{Contexto, público e objetivos}

\preencher{Descreva o público-alvo, o nível de escolaridade dos autores, a situação de uso e o escopo adotado. Deixe explícito que se trata de uma versão acadêmica para dois jogadores e informe o que ficou fora do escopo.}

O objetivo geral foi \preencher{escreva o objetivo em uma frase verificável}. Como objetivos específicos, buscou-se: \preencher{liste de três a cinco resultados concretos, como modelar regras, criar uma interface jogável, persistir partidas e construir uma árvore de estados}.

\section{Processo de desenvolvimento e entregas}

\preencher{Explique brevemente como o trabalho foi organizado: ciclos, responsáveis, critérios de conclusão, testes e fontes usadas para interpretar as regras. Em seguida, relate cada entrega com a mesma lógica: necessidade, implementação, decisão e evidência.}

\subsection{Primeira entrega: versão jogável e arquitetura MVC}

\textbf{O que foi feito.} \preencher{Resuma as entidades modeladas, as regras implementadas, o fluxo para dois jogadores, a interface, o menu e o salvamento.}

\textbf{Como foi feito.} \preencher{Explique a separação entre Model, View e Controller; a representação do tabuleiro como grafo; o estado completo da partida; e como as regras foram validadas. Cite classes ou módulos apenas quando isso ajudar o leitor.}

\textbf{Decisões e justificativas.} \preencher{Justifique MVC, estado centralizado, separação entre regra e interface e uso de uma configuração reproduzível. Relacione cada decisão a um benefício ou risco mitigado.}

\textbf{Resultado da entrega.} \preencher{Informe o que podia ser demonstrado ao fim desta etapa e quais limitações permaneceram.}

\subsection{Segunda entrega: jogadas completas e árvore de estados}

\textbf{O que foi feito.} \preencher{Descreva a geração e execução de jogadas legais, a árvore de estados, o retorno a estados anteriores, a eliminação de repetições, a persistência em JSON e as visualizações.}

\textbf{Como foi feito.} \preencher{Explique o fluxo Gerador de Jogadas $\rightarrow$ Executor $\rightarrow$ cópia do estado $\rightarrow$ comparação por identificador $\rightarrow$ novo nó. Informe que a expansão ocorre sob demanda.}

\textbf{Decisões e justificativas.} \preencher{Justifique por que cada nó representa uma jogada completa; por que estados são copiados; por que há um identificador estável; e por que alternativas são geradas sob demanda.}

\textbf{Resultado da entrega.} \preencher{Apresente evidências: navegação na árvore, salvamento/carregamento, exportação Graphviz e quantidade atualizada de testes aprovados.}

% Repita a subseção abaixo para entregas futuras, preservando o limite de páginas.
% \subsection{Terceira entrega: título}
% \textbf{O que foi feito.} ...
% \textbf{Como foi feito.} ...
% \textbf{Decisões e justificativas.} ...
% \textbf{Resultado da entrega.} ...

\section{Tecnologias e razões de escolha}

\begin{table}[ht]
  \centering
  \caption{Tecnologias empregadas no projeto.}
  \label{tab:tecnologias}
  \small
  \begin{tabularx}{\textwidth}{@{}p{0.19\textwidth}p{0.25\textwidth}Y@{}}
    \toprule
    \textbf{Tecnologia} & \textbf{Onde foi usada} & \textbf{Motivo da escolha} \\
    \midrule
    Python 3.13 & Regras, entidades, árvore e persistência & Legibilidade, biblioteca padrão abrangente e adequação ao contexto didático. \\
    Pygame 2.6.1 & Interface, entrada e ciclo principal & Permite criar uma interface 2D interativa mantendo controle explícito do ciclo do jogo. \\
    JSON & Salvamento da árvore e da partida & Formato textual, portável e compatível com dicionários e listas produzidos pelos serializadores. \\
    Graphviz & Exportação da árvore para DOT/SVG & Automatiza o posicionamento de grafos e produz diagramas vetoriais legíveis. \\
    \preencher{tecnologia} & \preencher{uso} & \preencher{justificativa vinculada ao requisito} \\
    \bottomrule
  \end{tabularx}
\end{table}

\preencher{Complemente a tabela com decisões que não cabem nela, como testes unitários, controle de versão e licenças. Evite dizer apenas que uma tecnologia é ``fácil''; compare a escolha com os requisitos do projeto.}

\section{Apresentação do software e resultados}

\preencher{Inclua uma ou duas capturas que mostrem o jogo e a árvore. Explique o fluxo visível, não apenas a aparência. Indique a versão avaliada, o ambiente de teste, a quantidade de testes executados e limitações conhecidas.}

% Exemplo de figura (adicione o arquivo à pasta figuras):
% \begin{figure}[ht]
%   \centering
%   \includegraphics[width=.88\textwidth]{figuras/interface-jogo.png}
%   \caption{Interface da partida e principais elementos de interação.}
%   \label{fig:interface}
% \end{figure}

\section{Considerações finais}

\preencher{Retome o objetivo e sintetize o aprendizado técnico e de projeto. Diferencie resultados observados de expectativas. Informe limitações e próximos passos, como função de avaliação, Minimax, poda alfa--beta e modo humano contra IA, somente se estiverem realmente planejados.}

% O MÁGICA 2026 exige de 3 a 5 páginas de conteúdo e admite uma página
% extra de referências.
% Confira o PDF final e ajuste o texto/figuras antes da submissão.

\bibliographystyle{sbc}
\bibliography{referencias}

\end{document}
